\section{BioseqDB indexing abstractions}
\label{sec:index-abstraction-layers}

Indexes are stored separately in BioseqDB~\ref{lst:types-of-storage-entities-in-bioseqdb} to all internal sequence fragmented tables~\ref{fig:storage-length-buckets} and they are separate entities
to the usual indexes in Postgres which provide generic access paths to the tables.

The index bytes are not stored in the fragment rows.
I keep them in one of three layouts.
A \emph{chunk} is one stored piece of those bytes.
One chunk must stay at or below $960$ MiB, which is $1$ GiB minus $64$ MiB, so the row still fits in Postgres.

\begin{itemize}
  \item \textbf{\texttt{CHUNKED\_PER\_FRAGMENT}}.
    One sequence can have many chunks.
    They are grouped by a \emph{sub-key}, an extra number chosen by the index.
    For \texttt{minimizers\_per\_fragment} that number is the fragment, so one fragment can fill several chunks.
    \texttt{minimizers\_per\_sequence} uses this same layout.
    Its sub-key is the number of one hash group, not a fragment number.
  \item \textbf{\texttt{INLINE\_PER\_SEQUENCE}}.
    One sequence has a single chunk.
    There is no sub-key and no second chunk.
    If the bytes exceed $960$ MiB, the write stops with an error.
    \texttt{symbol\_counts} uses this layout.
    It stores how many times each symbol occurs in that sequence.
  \item \textbf{\texttt{CHUNKED\_PER\_TABLE}}.
    The chunks belong to the whole user table.
    A build can visit every sequence.
    It is not maintained automatically when a row is inserted.
    No production index uses it yet.
    A test index uses it when tests are compiled in.
\end{itemize}

Figure~\ref{fig:index-storage-layouts} draws the three layouts under one user table and the fragment rows of one sequence.

The three layouts go through one storage interface, so a later index can pick a layout without a new file format.
The check is the C++ concept \texttt{is\_index\_chunked\_storage\_for}.
The fragment calls take a sub-key and a chunk number.
The inline calls do not take a sub-key or a chunk number.
The table calls take a chunk number, and a build can walk every sequence.
Postgres implements one class for each layout:
\begin{itemize}
  \item \texttt{fragment\_index\_storage} for many chunks under a sub-key
  \item \texttt{sequence\_index\_storage} for the single chunk
  \item \texttt{table\_index\_storage} for chunks of the whole table
\end{itemize}
These classes create the relations. A new index does not need a new table definition. The first use of its name creates the tables in the \texttt{bioseq\_internal} schema.

Each index name has its own family of tables.
\begin{itemize}
  \item A per-fragment or per-sequence relation is named \texttt{seqidx\_}, then a short name, then a hash, then the length bucket and the shard.
  \item \texttt{minimizers\_per\_fragment} is shortened to \texttt{mpf}. \texttt{minimizers\_per\_sequence} is shortened to \texttt{mps}. \texttt{symbol\_counts} is shortened to \texttt{sc}.
  \item The hash mixes the database, the tablespace, the user table's relfilenode, the sequence column, and the full index name. Two index names therefore do not share a relation.
  \item The bucket and the shard come from the sequence locator. They are the same length bucket and shard as that sequence's fragment rows.
\end{itemize}
\texttt{mps} does not add a new table shape. It uses the same per-fragment storage class as \texttt{mpf}, so the columns match, but the relations are a separate family.
The inline layout uses the same naming scheme. Its primary key is only the sequence row number, because that sequence has one chunk.
The per-fragment primary key also contains the sub-key and the chunk number.
The table-wide layout uses the prefix \texttt{seqidxtab\_} instead of \texttt{seqidx\_}. Its primary key is only the chunk number.

These relations do not use the columnar access method. They are created with \texttt{USING heap}. The user table is the columnar table. A minimizer generation relation is a heap table as well, and that one is created \texttt{UNLOGGED}. The generation check refuses the relation when it is not the heap access method.

\texttt{pg\_supported\_indexes} connects those classes to the registered indexes.
Listing~\ref{lst:mpf-index-storage-check} is the check used by minimizers per fragment.
The compiler rejects that index unless the storage provides the fragment-layout calls.
The listing leaves out the build functions.

\begin{BioseqCppSnippet}[caption={Minimizers per fragment require fragment-chunk storage.},
  label={lst:mpf-index-storage-check}]
template <class p_storage>
struct minimizers_per_fragment_index {
  static_assert(is_index_chunked_storage_for<
      p_storage, sequence_index_layout::fragment_chunks>);
  static constexpr sequence_index_layout index_type =
      sequence_index_layout::fragment_chunks;
  using storage_type = p_storage;
};
\end{BioseqCppSnippet}

\begin{figure}[htbp]
\centering
\begin{tikzpicture}[
  x=1cm,y=1cm,font=\footnotesize,
  lab/.style={align=center,font=\scriptsize,inner sep=1pt},
  arr/.style={line width=0.85pt,line join=miter,line cap=rect,
    -{Latex[length=2.4mm,width=1.7mm]}},
  stem/.style={line width=0.85pt,line join=miter,line cap=rect}
]
  \node[anchor=south west,font=\bfseries\footnotesize] at (0.85,12.22)
    {(a) User table};
  \draw[indexGreen,line width=1pt,fill=black!3]
    (0.85,9.35) rectangle (6.35,12.05);
  \node[lab,anchor=west] at (1.05,11.68) {id};
  \node[lab,anchor=west] at (2.45,11.68) {sequence locator};
  \draw[black!25] (0.85,11.42) -- (6.35,11.42);
  \fill[black!5] (1.05,10.72) rectangle (2.05,11.22);
  \draw[black!40] (1.05,10.72) rectangle (2.05,11.22);
  \node[lab] at (1.55,10.97) {2};
  \fill[black!5] (2.25,10.72) rectangle (6.15,11.22);
  \draw[black!40] (2.25,10.72) rectangle (6.15,11.22);
  \node[lab] at (4.20,10.97) {locator of sequence 2};
  \fill[black!5] (1.05,9.58) rectangle (2.05,10.28);
  \draw[black!40] (1.05,9.58) rectangle (2.05,10.28);
  \node[lab] at (1.55,9.93) {1};
  \fill[indexOrange!22] (2.25,9.58) rectangle (6.15,10.28);
  \draw[indexOrange,line width=1pt] (2.25,9.58) rectangle (6.15,10.28);
  \node[lab] at (4.20,9.93) {locator of sequence 1};

  % ---------- (b) fragments ----------
  \node[anchor=south west,font=\bfseries\footnotesize] at (7.15,12.22)
    {(b) Fragment rows of sequence 1};
  \draw[black!55,fill=black!3] (7.15,9.35) rectangle (15.05,12.05);
  \fill[black!5] (7.40,9.72) rectangle (9.55,11.55);
  \draw[black!40] (7.40,9.72) rectangle (9.55,11.55);
  \node[lab] at (8.475,10.64) {fragment 0};
  \fill[indexBlue!18] (9.80,9.72) rectangle (12.35,11.55);
  \draw[indexBlue,line width=1pt] (9.80,9.72) rectangle (12.35,11.55);
  \node[lab] at (11.075,10.64) {fragment 1};
  \fill[black!5] (12.60,9.72) rectangle (14.80,11.55);
  \draw[black!40] (12.60,9.72) rectangle (14.80,11.55);
  \node[lab] at (13.70,10.64) {fragment 2};

  % ---------- shared row geometry ----------
  % bar i top = 4.28 - i*0.56, height 0.46, center = 4.05 - i*0.56

  % ---------- (c) table chunks: all 5 rows ----------
  \draw[indexGreen!70,fill=indexGreen!5] (1.05,1.28) rectangle (4.85,5.62);
  \node[lab,font=\bfseries\scriptsize] at (2.95,5.22) {\texttt{CHUNKED\_PER\_TABLE}};
  \node[lab] at (2.95,4.72) {chunks of the whole table};
  \foreach \i/\name in {0/{chunk 0},1/{chunk 1},2/{chunk 2},3/{chunk 3},4/{chunk 4}} {
    \fill[indexGreen!20] (1.22,{4.28-\i*0.56}) rectangle (4.68,{3.82-\i*0.56});
    \draw[indexGreen] (1.22,{4.28-\i*0.56}) rectangle (4.68,{3.82-\i*0.56});
    \node[lab] at (2.95,{4.05-\i*0.56}) {\name};
  }

  % ---------- (d) inline: 5 rows, one belongs to sequence 1 ----------
  \draw[black!45,fill=indexOrange!4] (6.05,1.28) rectangle (9.85,5.62);
  \node[lab,font=\bfseries\scriptsize] at (7.95,5.22) {\texttt{INLINE\_PER\_SEQUENCE}};
  \node[lab] at (7.95,4.72) {one row for each sequence};
  \foreach \i/\name/\on in {0/{sequence 2}/0,1/{sequence 1}/1,2/{other sequence}/0,3/{other sequence}/0,4/{other sequence}/0} {
    \ifnum\on=1
      \fill[indexOrange!24] (6.22,{4.28-\i*0.56}) rectangle (9.68,{3.82-\i*0.56});
      \draw[indexOrange,line width=0.9pt] (6.22,{4.28-\i*0.56}) rectangle (9.68,{3.82-\i*0.56});
    \else
      \fill[black!4] (6.22,{4.28-\i*0.56}) rectangle (9.68,{3.82-\i*0.56});
      \draw[black!35] (6.22,{4.28-\i*0.56}) rectangle (9.68,{3.82-\i*0.56});
    \fi
    \node[lab] at (7.95,{4.05-\i*0.56}) {\name};
  }

  % ---------- (e) fragment chunks: 3 of 5 rows are fragment 1 ----------
  \draw[black!45,fill=indexBlue!4] (10.85,1.28) rectangle (15.05,5.62);
  \node[lab,font=\bfseries\scriptsize] at (12.95,5.22) {\texttt{CHUNKED\_PER\_FRAGMENT}};
  \node[lab] at (12.95,4.72) {chunks of fragment 1};
  \foreach \i/\name/\on in {
      0/{fragment 0}/0,
      1/{fragment 1, chunk 0}/1,
      2/{fragment 1, chunk 1}/1,
      3/{fragment 1, chunk 2}/1,
      4/{fragment 2}/0} {
    \ifnum\on=1
      \fill[indexBlue!18] (11.02,{4.28-\i*0.56}) rectangle (14.88,{3.82-\i*0.56});
      \draw[indexBlue] (11.02,{4.28-\i*0.56}) rectangle (14.88,{3.82-\i*0.56});
    \else
      \fill[black!4] (11.02,{4.28-\i*0.56}) rectangle (14.88,{3.82-\i*0.56});
      \draw[black!35] (11.02,{4.28-\i*0.56}) rectangle (14.88,{3.82-\i*0.56});
    \fi
    \node[lab] at (12.95,{4.05-\i*0.56}) {\name};
  }

  % Green: open brace on the whole user table, then one branch into every chunk row.
  % End caps stop short of the table border so they do not form an extra column.
  \draw[stem,indexGreen] (0.62,11.92) -- (0.42,11.92) -- (0.42,10.70)
    -- (0.28,10.70) -- (0.28,1.81);
  \draw[stem,indexGreen] (0.62,9.48) -- (0.42,9.48) -- (0.42,10.70);
  \foreach \i in {0,1,2,3,4} {
    \draw[arr,indexGreen] (0.28,{4.05-\i*0.56}) -- (1.22,{4.05-\i*0.56});
  }

  % Orange: locator cell -> the one sequence-1 row (bar 1, center 3.49).
  \draw[stem,indexOrange] (4.20,9.58) -- (4.20,8.55) -- (5.45,8.55) -- (5.45,3.49);
  \draw[arr,indexOrange] (5.45,3.49) -- (6.22,3.49);

  % Blue: leave fragment 1 itself, then split into its three chunk rows.
  \draw[stem,indexBlue] (11.075,9.72) -- (11.075,7.55) -- (10.22,7.55) -- (10.22,2.37);
  \foreach \i in {1,2,3} {
    \draw[arr,indexBlue] (10.22,{4.05-\i*0.56}) -- (11.02,{4.05-\i*0.56});
  }
\end{tikzpicture}
\caption[Three index storage layouts]{Each index is drawn as five rows.
The green brace marks the whole user table, and its line splits into every chunk row.
The orange arrow leaves the locator of sequence 1 and ends on that sequence's one row.
That row cannot grow into a second row.
The blue arrow leaves fragment 1 and splits into that fragment's three chunk rows.
Grey rows belong to other sequences or other fragments.
This drawing shows the fragment-number use of \texttt{CHUNKED\_PER\_FRAGMENT}.
The same layout can use another sub-key.
For \texttt{minimizers\_per\_sequence} that sub-key is the number of one hash group.}
\label{fig:index-storage-layouts}
\end{figure}


% The sequence store can fetch a small interval without decoding a whole
% chromosome. Seed lookup has a different question: at which positions does a
% query minimizer occur in the target? Reading every target symbol and rebuilding
% its minimizer table for each query would repeat much of the same work.
% A persistent index saves that table between executions.

% The useful object is a mapping from a seed key to its occurrences. A
% \emph{posting} is one stored occurrence, including a position and any source or
% strand information needed to interpret it. A \emph{posting list} contains all
% stored occurrences for one key. The key may be a minimizer hash, as in
% BioSeqDB and minimap2, or a spaced word, as in the CPU MMseqs2 prefilter.
% The seed-selection rule determines which occurrences the index stores.

% Reuse of reference indexes already appears in biological workflows. For
% example, BioWorkbench's RASflow workflow checks for an existing Bowtie2 index
% before preparing it again\autocite{Mondelli_2018}. MMseqs2's interactive
% search system also saves precomputed search structures to reduce setup for
% repeated requests\autocite{mmseqs2_desktop_2019}. BioSeqDB has to connect this
% reuse to the database's changing rows and to the search operators from the
% previous chapters.

% Four questions guide the design:
% \begin{itemize}
%   \item \textbf{What is stored?} Minimizer positions answer seed lookups;
%     symbol counts answer a different question about sequence composition.
%   \item \textbf{How is it grouped?} Grouping by target interval makes
%     construction local. Grouping by hash makes later reads selective.
%   \item \textbf{When is it usable?} The index must describe the same source
%     version and minimizer rules as the query expects.
%   \item \textbf{Who keeps its bytes alive?} A reader needs valid storage and
%     buffers until its last lookup, including when another transaction changes
%     the source.
% \end{itemize}

% An index is derived data: removing it does not remove the sequence. This also
% distinguishes it from the completed-query cache in
% Section~\ref{sec:bioseq-result-cache}. A minimizer index can serve many
% different query sequences; a cached result describes completed executions.

\section{BioseqDB supported indexing}
\label{sec:index-logical-layouts}

BioSeqDB supports \emph{minimizers per fragment} (MPF) and
\emph{minimizers per sequence} (MPS). MPF groups entries by source interval;
MPS groups them by hash across one logical sequence. Both preserve positions
needed to form anchors. Figure~\ref{fig:index-mpf-mps} shows the same six
occurrences in the two organizations. Here a minimizer position is its
last-symbol coordinate, using zero-based positions.

\begin{figure}[htbp]
\centering
\begin{tikzpicture}[x=1cm,y=1cm,font=\footnotesize]
  \node[anchor=west,font=\bfseries] at (0,6.05) {(a) MPF: keep each source interval together};
  \foreach \x/\f/\range in {2/0/{[0,100)},6/1/{[100,200)},10/2/{[200,300)}} {
    \draw[fill=black!3,draw=black!50] (\x-1.55,3.45) rectangle (\x+1.55,5.5);
    \node at (\x,5.15) {fragment \f: $\range$};
  }
  \node[index blue,minimum width=2.5cm] at (2,4.45) {$a:\ 17,65$};
  \node[index orange,minimum width=2.5cm] at (6,4.45) {$b:\ 120$};
  \node[index blue,minimum width=2.5cm] at (6,3.82) {$a:\ 180$};
  \node[index blue,minimum width=2.5cm] at (10,4.45) {$a:\ 230$};
  \node[index green,minimum width=2.5cm] at (10,3.82) {$c:\ 280$};
  \node[anchor=west,font=\bfseries] at (0,2.5) {(b) MPS: put equal hashes together across the sequence};
  \node[index blue,minimum width=1.0cm] at (1,1.65) {$a$};
  \foreach \x/\p in {3/17,4.6/65,6.2/180,7.8/230}
    \node[index blue,minimum width=1.45cm] at (\x,1.65) {$\p$};
  \node[index orange,minimum width=1.0cm] at (1,.85) {$b$};
  \node[index orange,minimum width=1.45cm] at (3,.85) {$120$};
  \node[index green,minimum width=1.0cm] at (1,.05) {$c$};
  \node[index green,minimum width=1.45cm] at (3,.05) {$280$};
  \foreach \y in {1.65,.85,.05}
    \draw[index arrow] (1.55,\y)--(2.22,\y);
  \node[anchor=west,align=left] at (9,1.0)
    {Query hash $a$ needs\\one sequence-wide list.};
\end{tikzpicture}
\caption{Changing the grouping changes which stored units a query needs.
Only the minimizer hashes $a,b,c$ are shown. All positions are shown in
global sequence coordinates to make the two panels comparable; MPF storage
can encode them relative to a fragment and rebase them when read. The
100-symbol intervals are illustrative.}
\label{fig:index-mpf-mps}
\end{figure}

\subsection{MPF follows source intervals}
\label{sec:index-mpf}

An MPF lookup has a hash table or ordered hash controls and a separate array
of positions. A \emph{control} records the hash and the location and length
of its posting list. Local construction avoids sorting the complete logical
sequence at once. A target query can read or assemble the necessary fragment
objects, then use their positions to construct anchors.

The ordinary MPF builder uses fixed 32,768-symbol cores. A \emph{core} is the
interval whose occurrences a construction task owns. A \emph{generation}
is a separately identified index build, which becomes usable only after
completion and publication. The optional grouped generation builder groups
positions according to physical source fragment
capacity, while collecting them through cores of at most 32,768 symbols.
These construction units should not be confused with a PostgreSQL page or
a compressed column chunk. Both builders need symbols around their boundaries
so that a minimizer near an edge is chosen from a complete window.
Section~\ref{sec:index-construction} explains this overlap.

The grouped MPF layout stores hashes, list offsets, counts, and positions in
separate arrays, also called a \emph{structure of arrays} (SoA). This gives
the reader explicit ranges to check and
allows a prepared lookup to use arrays without recreating a C++ hash table.
An additional read-only index file (Section~\ref{sec:index-mpf5-sidecar})
contains a sequence-wide hash-to-positions section built during
publication. It uses membership bits and compact list records, explained in
Section~\ref{sec:index-direct-lookup}. The lookup setting chooses this section
or the fragment controls in the same file. Building the section moves the
merge of equal hashes across fragments into index construction.
Grouped layouts and their readers are optional, rather than implied merely
by selecting the MPF index name.

\subsection{MPS follows hash ranges}
\label{sec:index-mps}

MPS combines the sequence's minimizers into hash order. Its directory maps
hash ranges to stored data. A query can search this directory before
loading positions. The current selective layout separates three streams:
\begin{itemize}
  \item the \textbf{directory} identifies the hash ranges and their controls;
  \item \textbf{controls} give hashes, exact occurrence counts, and posting
    locations;
  \item \textbf{postings} give the sequence positions and strands.
\end{itemize}
All indexed occurrences remain stored. A query's occurrence threshold controls
which posting lists are read; it does not remove them permanently during
construction.

MPS spends more work on sorting and merging so later queries can skip
unrelated hash ranges. MPF keeps its base representation closer to source
intervals. The current default active minimizer index is MPF. MPS is an
explicit choice for workloads that justify a sequence-wide build. The
registered persistent minimizer profile uses $k=15,w=10$; the selective MPS
stream route additionally requires the supported DNA4 and hash-layout
identity. Not every minimizer setting for a lookup built during the query
has a corresponding persistent reader.

\subsection{A smaller index answers symbol counts}
\label{sec:index-symbol-counts}

Not every sequence index stores seed positions. The \texttt{symbol\_counts}
index keeps a histogram for one sequence: one 64-bit counter for each encoded
symbol rank. For DNA4, the counters are $A,C,G,T$. The sequence
\texttt{AACGTA} therefore gives $(3,1,1,1)$, occupying 32 counter bytes before
row and value metadata.

\begin{figure}[htbp]
\centering
\begin{tikzpicture}[x=1cm,y=.7cm,font=\footnotesize]
  \node[anchor=west,font=\bfseries] at (0,5.6) {One sequence};
  \foreach \x/\s in {1/A,2/A,3/C,4/G,5/T,6/A}
    \node[index blue] at (\x,4.7) {\texttt{\s}};
  \node[anchor=west,font=\bfseries] at (8,5.6) {Four counters};
  \foreach \x/\s/\n in {8.5/A/3,9.7/C/1,10.9/G/1,12.1/T/1} {
    \draw[fill=indexGreen!18,draw=indexGreen] (\x-.38,.7) rectangle (\x+.38,.7+\n);
    \node[above] at (\x,.7+\n) {$\n$};
    \node[below] at (\x,.7) {\texttt{\s}};
  }
  \draw[index arrow] (6.5,3.6)--(7.75,3.6);
  \node[anchor=west,align=left] at (0,1.6)
    {To count A, read A's counter.\\Seed lookup: positions are still needed.};
\end{tikzpicture}
\caption{A histogram preserves sequence composition but loses position and
order. It can answer a count directly, but cannot locate an alignment seed.}
\label{fig:index-symbol-counts}
\end{figure}

\begin{lstlisting}[style=bioseq-index-sql,
  caption={Request a symbol count for each stored sequence.},
  label={lst:index-symbol-counts}]
SELECT id, bioseq.symbol_count(seq, 'A') AS a_count
FROM reads;
\end{lstlisting}
For a supported scan projection on a typed sequence column and a constant
letter, the executor can read the current row's saved histogram and supply
the selected counter. If the eligible histogram is unavailable, it counts
the lazy sequence's symbols. This index avoids a full sequence read for such
projections. It is not an implemented rule for discarding alignment
candidates, and it does not preserve enough information to answer position
queries. The histogram is stored in an internal relation keyed by the
parent's internal row number.
